% ---------------------------------------------------------------------------
\section*{Decision summary}
\addcontentsline{toc}{section}{Decision summary}
\begin{enumerate}[leftmargin=*,itemsep=2pt]
\item \textbf{No universal cavitation-specific maximum; conditional answers.} ``Highest pressure or jet speed'' has
an application-specific finite answer only after fixing the activation mechanism, the observable (seven are distinguished, O1--O7) and,
for speeds, the participating mass. This report fixes them for a declared benchmark (one fully activated
PFP or PFH droplet in water near a rigid receiver). It connects an assumed activation state to
exergy, expansion, reduced collapse dynamics and an internal directed microjet. Laser deposition
and receiver loading are respectively an adopted preparation and an order-of-magnitude wave estimate;
they are not resolved end-to-end device calculations.
\item \textbf{Conditional thermodynamic ceilings.} Under the stated dead-state and no-external-work
assumptions, only the exergy $\Phi_{\mathrm a}$ of the
activated droplet can become work, cavity energy, jet kinetic energy or fracture work. For homogeneous
activation (CNT kinetic limit $T_{\mathrm k}=\PFPTkSix$\,K PFP, $\PFHTkSix$\,K PFH, $\approx0.92T_{\mathrm c}$)
$\Phi_{\mathrm a}/Q=\PFPetaSix$ (PFP) and $\PFHetaSix$ (PFH): at most about one seventh of the activation heat is
mechanically available. Under H0, H2 and the equilibrium/no-overshoot hypothesis H5, the cavity
pressure cannot exceed $\PFPpsrcSp$\,MPa (PFP) or $\PFHpsrcSp$\,MPa (PFH) absolute during expansion
(Prop.~\ref{prop:psrc}); its extension to finite-rate HKS content is unproved. Under the stable-EOS
and reference-state hypotheses, the cavity radius cannot exceed $\PFPRratioSix\,a_0$ (PFP; Eq.~\ref{eq:vmax}),
which the inviscid dynamics attains to within $\PFPRsharpDyn$\%.
\item \textbf{Collapse is controlled by one unmeasured parameter.} Spherical real-fluid simulations (216 cases) show
that the PFC-vapour condensation coefficient $\alpha_{\mathrm c}$ decides whether the collapse is cushioned
($\alpha_{\mathrm c}\le0.1$: wall speeds $\sim$12--14\,m/s) or reaches the compressible regime
(6\,$\mu$m water samples: PFP at $\alpha_{\mathrm c}=0.5$, PFH already at $0.3$).
Thermal refinement leaves this classification unchanged but the transitional PFP
$\alpha_{\mathrm c}=0.3$ speed changes $\kmRefineMaxUncensoredDeltaU$\% from 160 to 320 nodes;
all six refined cases carry an ideal-vapour warning.
Gels arrest the expansion of small PFP droplets before full vaporisation (3 and 6\,$\mu$m at 20\,kPa, 3\,$\mu$m at
5\,kPa), removing the violent collapse in this model.
\item \textbf{Directed microjets (benchmarked boundary-integral solver).} For sampled stand-offs of a 6\,$\mu$m droplet at
$1.25\le\gamma\le4$ the predicted microjet speeds are $\PFPenvUmin$--$\PFPenvUmax$\,m/s (PFP) and
$\PFHenvUmin$--$\PFHenvUmax$\,m/s (PFH) within the Mach-0.3 validity screen, carrying $\PFPenvEfmin$--$\PFPenvEfmax$\% of
$\Phi_{\mathrm a}$ (PFP). With slow condensation the largest sampled speeds are $\PFPenvSlowUmax$\,m/s (PFP) and
$\PFHenvSlowUmax$\,m/s (PFH) on the original $N=64$ grid; selected $N=96$ controls give
$\PFPbimRefinedSlowSample$ and $\PFHbimRefinedSlowSample$\,m/s.
Three requested unchanged-cap spatial attempts abort; four separate smaller-cap $N=96$
follow-ups succeed, while the finer $N=128$ attempt aborts before impact.
Mass/profile-energy sensitivity is larger than speed sensitivity (Section~\ref{sec:bim_refine}).
Fast-condensation samples straddle the adopted incompressible-model
screen $0.3c_{\mathrm w}=\Ulim$\,m/s inside the design range (interpolated $\gamma\approx\PFPcrossIsog$ PFP,
$\PFHcrossIsog$ PFH). This is a \emph{reduced-model validity crossing}, not a demonstrated physical
lower bound or global maximum. Out-of-screen incompressible estimates are $\PFPenvOverMach$ (PFP)
and $\PFHenvOverMach$\,m/s (PFH) at $\gamma=4$. The local
liquid--liquid hammer pressure at jet impact reaches $\sim\PFPenvPhMax$\,MPa (PFP) and $\sim\PFHenvPhMax$\,MPa (PFH) for
$\sim$1\,ns in the adopted mixed-contact estimate (nonlinear jet-side Tait, linear opposing
impedance); the receiver-wave estimate is of order
$\PFPenvPrecMin$--$\PFHenvPrecMax$\,MPa, without a resolved propagation or sensor calculation.
\item \textbf{Audit corrections to prior work.} Limits E02 (budget must be exergy, not undefined availability),
E19 (empty-bubble radius bound is not a bound for a vapour cavity), E10 (linear water hammer underestimates by
up to $2.1\times$ at 850\,m/s within the adopted Tait extrapolation) and E25 (the sourced interpolated
heat capacity is replaced by the EOS-consistent value; a data/model discrepancy remains).
Homogeneous CNT activation is superheat to $\approx0.92T_{\mathrm c}$,
not boiling at 303\,K. The finite-PFC study's CNT gate is evaluated and its seeds have negligible
homogeneous-nucleation probability under the adopted capillarity/rate assumptions.
\item \textbf{Remaining gap (honestly unresolved).} The actual device maximum needs: the absorbed energy and
heating distribution, the nucleation route (heterogeneous sites lower $T_{\mathrm k}$ and $\Phi_{\mathrm a}$), the
condensation coefficient, the geometry/stand-off and the receiver/sensor. A seven-item measurement program
(Section~\ref{sec:program}) targets each gap; compressible focusing, gel fracture and receiver/sensor
models must also be completed where their reduced approximations fail. Optical breakdown and
near-contact needle jets are different mechanisms outside this benchmark.
\end{enumerate}
